\chapter{1904:Sir William Ramsay}
\label{chap:chemistry1904}

The Nobel Prize in Chemistry for the year 1904 was awarded to Sir William Ramsay.

\textbf{Motivation:}
in recognition of his services in the discovery of the inert gaseous elements in air, and his determination of their place in the periodic system

\section{Sir William Ramsay}

\subsection*{Early Life and Education}

Sir William Ramsay was born on 1852-10-02 in Glasgow, Scotland.
Sir William Ramsay was a Scottish chemist who discovered the noble gases and received the Nobel Prize in Chemistry in 1904 "in recognition of his services in the discovery of the inert gaseous elements in air" along with his collaborator, John William Strutt, 3rd Baron Rayleigh, who received the Nobel Prize in Physics that same year for their discovery of argon. After the two men identified argon, Ramsay investigated other atmospheric gases. His work in isolating argon, helium, neon, krypton, and xenon led to the development of a new section of the periodic table.

\subsection*{Career and Major Research}

Ramsay was professor of chemistry at University College London for much of his career. Working with Lord Rayleigh, he isolated argon from air in 1894, and in 1895 he found helium in the mineral cleveite. With Morris Travers he then discovered neon, krypton, and xenon in liquid air, and in 1903 he identified helium produced by radioactive decay, completing a new group of the periodic table.

\subsection*{Nobel Prize-winning Work}

The work for which Sir William Ramsay was awarded the Nobel Prize in Chemistry in 1904 was described by the Nobel Committee as: "in recognition of his services in the discovery of the inert gaseous elements in air, and his determination of their place in the periodic system".

This research fundamentally advanced the field by discovering the noble gases (argon, helium, neon, krypton, xenon) and determining their place in the periodic system, completing a major gap in the periodic table.

\subsection*{Significance and Impact}

The impact of Sir William Ramsay's work extends far beyond the specific achievement recognized by the Nobel Prize.
The discovery of the noble gases completed the periodic table and led to the development of a new group of elements with unique properties, influencing both theoretical chemistry and practical applications in lighting and cryogenics.

\subsection*{Later Career and Honors}

Sir William Ramsay continued his scientific work until 1916-07-23, passing away in High Wycombe, United Kingdom.
He received numerous honors including membership in major scientific academies and societies.
Sir William Ramsay was knighted in 1902 and served as president of the Chemical Society and the British Association for the Advancement of Science.

\subsection*{Legacy}

The legacy of Sir William Ramsay endures through the lasting impact of his scientific contributions.
The noble gases, once considered chemically inert, now form compounds that have expanded the boundaries of chemistry. Sir William Ramsay's discovery completed a fundamental gap in the periodic system.

\subsection*{Selected Publications}

"The Gases of the Atmosphere" (1896)
"On a New Gas, Argon" (1895)
"Helium, a Gaseous Constituent of Certain Minerals" (1895)

\clearpage

\section*{Chapter Summary}

Sir William Ramsay was honored for the discovery of the inert gaseous elements of air, including argon, helium, neon, krypton, and xenon, completing a new group of the periodic table.
